\section{Q-Learning under Constant Valuations}
\label{sec:exp1_results}

Experiment~1 deploys Q-learning agents with constant valuations ($v=1$) in a $2^{10-1}$ Resolution~V half-fraction spanning 10 factors and 1{,}024 observations. Under independent private values with constant valuations, revenue equivalence \citep{Vickrey1961} predicts identical expected revenue across auction formats. Despite this theoretical equivalence, average revenues in the final 1{,}000 episodes range from approximately 0.3 to 1.0, revealing substantial variability across configurations of learning parameters and market structure. The factorial analysis decomposes this variability into systematic effects attributable to individual factors and their interactions.

\begin{figure}[H]
  \centering
  \includegraphics[width=\textwidth]{figures/e1_trace}
  \caption{Representative learning trajectory for a single trial of Experiment~1 (first-price auction, 2~bidders, Boltzmann exploration, 10{,}000 episodes). Per-bidder mean bids (200-episode rolling mean) converge toward the BNE prediction over the course of training.}
  \label{fig:e1_trace}
\end{figure}

\subsection{Revenue}

Number of bidders is the dominant factor for revenue ($|t| = \ExpOneRevFactorTOne$), followed by the discount factor $\gamma$ ($|t| = \ExpOneRevFactorTTwo$) and update mode ($|t| = \ExpOneRevFactorTThree$). Table~\ref{tab:exp1_ranked_rev} ranks significant effects by absolute $t$-statistic. Moving from two to four bidders increases average revenue by \ExpOneRevTopPctEffect\% of the grand mean, an effect \ExpOneRevTopVsAuctionRatio{} times larger than that of auction format. Auction format, by contrast, does not exert a statistically significant main effect on revenue ($|t| = \ExpOneRevAuctionT$, $p \ExpOneRevAuctionPFmt$). First-price and second-price auctions produce nearly identical average revenues in the final 1{,}000 episodes (FPA mean \ExpOneRevMeanFPA{} vs.\ SPA mean \ExpOneRevMeanSPA). The main effects plot (Figure~\ref{fig:e1_main_rev}) confirms these findings.

\input{tables/exp1_ranked_rev}

\begin{figure}[H]
  \centering
  \includegraphics[width=0.7\textwidth]{figures/e1_main_rev}
  \caption{Experiment~1: Main effects plot for average revenue. Number of bidders and discount factor dominate; auction format has a negligible main effect.}
  \label{fig:e1_main_rev}
\end{figure}

Although auction format has no significant main effect on final-episode revenue, it does interact significantly with other factors. The auction format $\times$ number of bidders interaction ($|t| = \ExpOneRevAuctionxNbidAbsT$) indicates that the auction format gap is larger in thin markets than in thick markets. The auction format $\times$ discount factor interaction ($|t| = \ExpOneRevAuctionxGammaAbsT$) shows that the auction format effect varies with the discount factor. Figure~\ref{fig:e1_int_rev} displays these top interaction effects.

\begin{figure}[H]
  \centering
  \includegraphics[width=0.7\textwidth]{figures/e1_int_rev}
  \caption{Experiment~1: Interaction plot for average revenue, showing the top two-way interactions. Non-parallel lines indicate that the effect of one factor depends on the level of another.}
  \label{fig:e1_int_rev}
\end{figure}

\subsection{Price Volatility}

Table~\ref{tab:exp1_ranked_vol} ranks the significant effects for price volatility. Market thickness, reserve price, discount factor, and update mode are the primary drivers (Figure~\ref{fig:e1_main_vol}). Neither auction format nor exploration strategy has a significant direct effect.

\input{tables/exp1_ranked_vol}

\begin{figure}[H]
  \centering
  \includegraphics[width=0.7\textwidth]{figures/e1_main_vol}
  \caption{Experiment~1: Main effects plot for price volatility. Number of bidders, reserve price, and discount factor are the primary drivers.}
  \label{fig:e1_main_vol}
\end{figure}

Model adequacy diagnostics (Appendix~\ref{sec:appendix_robustness}) confirm that the linear model captures the dominant effects.\footnote{$R^2$ values range from \ExpOneRsqMin{} to \ExpOneRsqMax{} across response variables. LightGBM cross-validated $R^2$ modestly exceeds OLS for revenue (\ExpOneRevLGBMRsq{} vs.\ \ExpOneRevRsq), and the gap exceeds the 0.05 suggestive threshold for several responses, indicating detectable nonlinear structure; however, effect rankings are unchanged under the nonparametric model. Of the \ExpOneNRawSig{} effects significant at $p < 0.05$, \ExpOneNBHSig{} (\ExpOneBHPct\%) survive Benjamini--Hochberg correction.}

\subsection{Summary}

The number of bidders is the strongest determinant of seller outcomes, followed by the discount factor and update mode. Auction format does not significantly affect end-state revenue but interacts with competitive pressure and temporal discounting. This null result for auction format provides a baseline. Subsequent experiments show that the importance of format is context-dependent, emerging under affiliated valuations and budget-constrained pacing. Price volatility is governed by market thickness, reserve prices, and the discount factor. Global sensitivity analysis independently confirms these rankings, with the number of bidders achieving the highest total-order Sobol' index for revenue ($S_T = 0.24$) and reserve price dominating the no-sale rate ($S_T = 0.27$), while auction format is negligible for revenue ($S_T = 0.045$). Cross-method concordance is high ($\rho = 0.87$), indicating that all seven independent sensitivity methods agree on the factor hierarchy (Appendix~\ref{sec:appendix_sensitivity}). Having established that auction format is revenue-neutral under constant valuations with Q-learning, Section~\ref{sec:exp2_results} tests whether this neutrality persists when valuations are interdependent, as \citet{MilgromWeber1982} predict that affiliated values should break revenue equivalence in favour of second-price auctions.
